\begin{helper}[FiberAndDegreeSameColorLiftedIntersectionNegligible]
\label{helper:FiberAndDegreeSameColorLiftedIntersectionNegligible}
Let \(\beta=1/2\), let
\[
m(n)=\noderef{TwoBiteNaturalReducedVertexCount}(n),
\qquad
p(n)=\noderef{TwoBiteEdgeProbability}(\beta,n),
\]
and take probabilities with respect to
\[
\noderef{TwoBiteEventProbability}(n,m(n),p(n)).
\]
The following event has \(\noderef{NegligibleProbability}\): there is a
distinct red pair or a distinct blue pair whose lifted-neighborhood
intersection is larger than \(100(\log n)^3\), and simultaneously every
same-color base codegree is at most \(100\log n\).  Equivalently, the
negligible event is
\[
\begin{aligned}
&\Bigl[
\exists r\ne s,
\left|
\noderef{TwoBiteLiftedNeighborhood}(\omega,\operatorname{inl}(r))
\cap
\noderef{TwoBiteLiftedNeighborhood}(\omega,\operatorname{inl}(s))
\right|>100(\log n)^3\\
&\quad\text{or}\quad
\exists b\ne c,
\left|
\noderef{TwoBiteLiftedNeighborhood}(\omega,\operatorname{inr}(b))
\cap
\noderef{TwoBiteLiftedNeighborhood}(\omega,\operatorname{inr}(c))
\right|>100(\log n)^3
\Bigr]\\
&\quad\cap
\Bigl[
\forall r\ne s,
|\{t:\noderef{TwoBiteRedGraph}(\omega).Adj(r,t)\land
\noderef{TwoBiteRedGraph}(\omega).Adj(s,t)\}|\le100\log n,\\
&\qquad\quad
\forall b\ne c,
|\{t:\noderef{TwoBiteBlueGraph}(\omega).Adj(b,t)\land
\noderef{TwoBiteBlueGraph}(\omega).Adj(c,t)\}|\le100\log n
\Bigr].
\end{aligned}
\]
\end{helper}

\begin{proof}
Write
\[
M=m(n),\qquad p=p(n),\qquad L=\log n.
\]
Let \(E_n\) be the event in the statement.  Let \(B_n\) be the sharper
same-color base-codegree event
\[
\begin{aligned}
B_n=\{&
\forall r\ne s,
|\{t:\noderef{TwoBiteRedGraph}(\omega).Adj(r,t)\land
\noderef{TwoBiteRedGraph}(\omega).Adj(s,t)\}|\le L,\\
&\forall b\ne c,
|\{t:\noderef{TwoBiteBlueGraph}(\omega).Adj(b,t)\land
\noderef{TwoBiteBlueGraph}(\omega).Adj(c,t)\}|\le L\}.
\end{aligned}
\]
Pointwise,
\[
E_n\subseteq (E_n\cap B_n)\cup B_n^c.
\tag{1}
\]
Thus it suffices to prove that both probabilities on the right tend to
\(0\).

First prove \(\Pr(B_n^c)\to0\).  Discard finitely many \(n\) so that
\[
L\ge1,\qquad M=\left\lfloor \frac n{L^2}\right\rfloor,\qquad
M\le n,\qquad 0\le p\le1.
\tag{2}
\]
The formula for \(M\) follows from \(\noderef{TwoBiteNaturalReducedVertexCount}\),
\(\noderef{TwoBiteReducedVertexCount}\), and \(\noderef{TwoBiteStretch}\);
the probability bounds on \(p\) are eventual by
\(\noderef{OppositeProjectionEdgeProbBounds}\).  Fix distinct red vertices
\(r,s\).  By \(\noderef{TwoBiteSample}\), an outcome is a red graph, a blue
graph, and an injection.  For any event depending only on the red graph, the
finite probability sum in \(\noderef{TwoBiteEventProbability}\), with the
factorization of \(\noderef{TwoBiteSampleWeight}\), separates into a red
\(\noderef{GnpGraphWeight}\) sum, a blue graph sum, and an injection sum.
The blue graph sum is \(1\) by \(\noderef{GnpGraphWeightSumOne}\), and the
injection sum is \(1\) by the definition of \(\noderef{UniformInjectionWeight}\)
once \(n\le M^2\), which holds eventually by
\(\noderef{TwoBiteNaturalTotalMassOneEventually}\).  Hence the red marginal
is the \(G(M,p)\) product law.  The same argument gives the blue marginal.

Under this product law, the red common-neighbor count of \(r,s\) is a binomial
variable with at most \(M\) trials and success probability \(p^2\).  Its mean
satisfies
\[
\mu\le Mp^2
\le \frac n{L^2}\cdot \frac14\frac Ln
=\frac1{4L}.
\tag{3}
\]
Let \(T=\lceil L\rceil\).  If the common-neighbor count is \(>L\), then it is
at least \(T\).  A direct union bound over \(T\)-subsets of possible common
neighbors gives
\[
\Pr(Y\ge T)\le \binom MT(p^2)^T
\le \left(\frac{e\mu}{T}\right)^T.
\tag{4}
\]
Using \(T\ge L\) and (3), for large \(n\),
\[
\Pr(Y>L)\le
\left(\frac{e}{4L^2}\right)^L
\le \exp\!\left(-\frac12L\log L\right).
\tag{5}
\]
There are fewer than \(M^2\le n^2=\exp(2L)\) ordered red pairs and the same
number of ordered blue pairs.  The union bound, (5), and the domination of
\(L\log L\) over \(L\) prove \(\Pr(B_n^c)\to0\).

It remains to bound \(E_n\cap B_n\).  Fix base graphs satisfying \(B_n\).
For a distinct red pair \(r,s\), membership in both lifted red neighborhoods
means exactly that the injected product cell has red coordinate in
\[
N_{\noderef{TwoBiteRedGraph}(\omega)}(r)\cap
N_{\noderef{TwoBiteRedGraph}(\omega)}(s),
\]
by the definitions of \(\noderef{TwoBiteLiftedNeighborhood}\),
\(\noderef{TwoBiteRedProjection}\), and \(\noderef{TwoBiteEmbedding}\).
Thus the marked product cells are
\[
(N_{G_R}(r)\cap N_{G_R}(s))\times\operatorname{Fin}(M),
\]
and \(B_n\) gives at most \(LM\) marked cells.

The rounded-size definition gives, after discarding finitely many more \(n\),
\[
n\le M^2,\qquad n\le 2L^2M.
\tag{6}
\]
The first inequality is again supplied by
\(\noderef{TwoBiteNaturalTotalMassOneEventually}\); the second is the
standard floor consequence of
\(M=\lfloor n/L^2\rfloor\) for large \(n\).  Therefore all hypotheses of
\(\noderef{FiberAndDegreeSameColorLiftedIntersectionFixedPairInjectionTail}\)
hold for this fixed red pair.  The conditional probability, over the
injection coordinate, that this pair has lifted intersection larger than
\(100L^3\) is at most \(\exp(-2L^3)\).

The blue case is identical.  Conditioning on the blue base graph, the marked
cells for a distinct blue pair \(b,c\) are
\[
\operatorname{Fin}(M)\times (N_{G_B}(b)\cap N_{G_B}(c)),
\]
using \(\noderef{TwoBiteBlueProjection}\), and \(B_n\) again gives at most
\(LM\) marked cells.  Hence the same fixed-pair tail bound applies.

There are fewer than \(M^2\) ordered red pairs and fewer than \(M^2\) ordered
blue pairs, so the conditional probability of \(E_n\cap B_n\), given any
base graphs satisfying \(B_n\), is at most
\[
2M^2\exp(-2L^3).
\tag{7}
\]
The sample weights are nonnegative for all sufficiently large \(n\) by
\(\noderef{TwoBiteSampleWeightNonnegative}\) and
\(\noderef{OppositeProjectionEdgeProbBounds}\), so the same finite-sum union
bound applies after summing over base graphs.  Thus the unconditional
probability of \(E_n\cap B_n\) is bounded by the right side of (7).
By \(\noderef{OppositeProjectionAsymptoticBound}\) with \(c=2\),
\[
M(n)^2\exp(-2(\log n)^3)\to0.
\]
Combining this with \(\Pr(B_n^c)\to0\) and the decomposition (1) proves that
the probability of \(E_n\) tends to \(0\), which is exactly the asserted
\(\noderef{NegligibleProbability}\).
\end{proof}
